% TOP_PROBLEM: {"id": 30007332, "title": "Uniqueness of admissible periodic hard-sphere Boltzmann evolution", "category_id": 9, "category": {"id": 9, "name": "pde", "display_name": "Partial Differential Equations", "description": "PDEs and their applications in physics and geometry.", "slug": "pde", "order_index": 9, "created_at": "2026-07-31T15:26:25.671Z"}, "status": "open"}
% FIRST_POSTED: 2026-10-06
% ATTEMPT_STATUS: solved
% OPENAI_CLAIM: {"schema_version":1,"source_repo":"https://github.com/openai/math","source_commit":"adc7f1241b42e322a6451854ab7e4b4c146bf78a","release_date":"2026-10-06","match":"stronger","resolution":"disproved","independently_reviewed":false,"second_pass":"agreed","families":[{"family":"363","title":"Nonuniqueness with local conservation for the hard-sphere Boltzmann equation","subject":"Partial differential equations","manuscripts":[{"title":"Nonuniqueness for the periodic hard-sphere Boltzmann equation","pdf_path":"preprints/Nonuniqueness-for-the-periodic-hard-sphere-Boltzmann-equation-September-23-2026/paper.pdf","readme_path":"preprints/Nonuniqueness-for-the-periodic-hard-sphere-Boltzmann-equation-September-23-2026/README.md","date":"2026-09-23","theorem_ref":"preprints/Nonuniqueness-for-the-periodic-hard-sphere-Boltzmann-equation-September-23-2026/build/sections/01-introduction.tex:55 (thm:main)"},{"title":"Nonuniqueness with local conservation for the hard-sphere Boltzmann equation","pdf_path":"preprints/Nonuniqueness-with-local-conservation-for-the-hard-sphere-Boltzmann-equation-October-5-2026/paper.pdf","readme_path":"preprints/Nonuniqueness-with-local-conservation-for-the-hard-sphere-Boltzmann-equation-October-5-2026/README.md","date":"2026-10-05","theorem_ref":"preprints/Nonuniqueness-with-local-conservation-for-the-hard-sphere-Boltzmann-equation-October-5-2026/build/paper/01-introduction.tex:160 (thm:main)"}],"lean":{"doc_path":"lean/docs/363.md","comparators":["lean/ComparatorChallenges/BoltzmannNonuniqueness.json","lean/ComparatorChallenges/BoltzmannNonuniqueness.lean"],"declarations":[],"covers_main_theorem":false},"reasoning_trace":null}]}
\documentclass[11pt]{article}
\usepackage[margin=25mm]{geometry}
\usepackage{fontspec}
\setmainfont{Latin Modern Roman}
\usepackage{amsmath,amssymb,mathtools}
\usepackage{unicode-math}
\setmathfont{Latin Modern Math}
\usepackage{enumitem}
\usepackage{xurl,hyperref,bookmark}
\hypersetup{colorlinks=true,urlcolor=blue}
\setcounter{secnumdepth}{0}
\setlength{\parindent}{0pt}
\setlength{\parskip}{5pt}
\setlength{\emergencystretch}{3em}
\begin{document}
\section*{Uniqueness of admissible periodic hard-sphere Boltzmann evolution}

\subsection{OpenAI's claim}
\textit{Periodic Boltzmann nonuniqueness}: The manuscript claims two global renormalized hard-sphere Boltzmann solutions on $\mathbb T^3\times\mathbb R^3$ with the same initial density; they satisfy the stated local-mass, total-momentum, energy, and entropy conditions, with strong $L^1$ continuity and integrable gains and losses on an initial interval.

\textit{Boltzmann nonuniqueness with local conservation}: The refinement claims two distinct global entropy solutions for one periodic hard-sphere initial datum while preserving mass, momentum, and kinetic energy locally; polynomially weighted gain and loss integrability and strong $L^1$ continuity hold on bounded time intervals.

\subsection{Outline of the argument}
\textit{Periodic Boltzmann nonuniqueness}: Two cutoff approximation families separate a dormant hot component from one activated by a seed of vanishing mass. Velocity-weight caps provide uniform integrability and permit removal of the annular and collision cutoffs. A rescaled branching argument prevents an apparent cap hit from disappearing into negligible mass, proving that the admissible limits differ despite sharing the same initial datum.

\textit{Boltzmann nonuniqueness with local conservation}: Kernel estimates and carefully arranged cold jets select short-time hot-component approximations. A global continuation bound and weighted compactness preserve collision terms and each local conservation law. Dilation and branching produce a nonzero rescaled limit with a finite horizon, which separates the two limiting solutions rather than merely their approximations.

\subsection{Scope relative to this problem}
\textit{Periodic Boltzmann nonuniqueness}: The announced pair shares one bounded-velocity datum satisfying the required mass, energy and entropy bounds. Both evolutions meet the renormalized admissibility and initial-interval integrability conditions, so their distinctness directly contradicts the uniqueness statement.

\textit{Boltzmann nonuniqueness with local conservation}: The announced pair shares one bounded-velocity datum satisfying the required mass, energy and entropy bounds. Both evolutions meet the renormalized admissibility and initial-interval integrability conditions, so their distinctness directly contradicts the uniqueness statement.

The claimed counterexample lies within the question’s hypotheses and contradicts its universal conclusion.

\subsection{Formal verification}
Family 363 describes Lean coverage for the initial periodic nonuniqueness result with its stated conservation and integrability properties; the stronger local-conservation refinement is separate.

\textit{Periodic Boltzmann nonuniqueness}: The manuscript links \texttt{BoltzmannNonuniqueness.lean}.

The documented selection includes this manuscript’s main assertion; that assertion retains the restrictions stated above.

\textit{Boltzmann nonuniqueness with local conservation}: No manuscript-specific Comparator is listed in the pinned inventory.

The documented selection does not cover the entire manuscript assertion relevant here.

This site has not run the Lean development. These links describe OpenAI’s selected formal statements and should be read with their stated scope.

\subsection{Status caveat}
OpenAI reports that an internal model produced the manuscripts. Their arguments have not been independently verified as proofs by this site. Any portion outside the documented formal scope may contain mistakes, and the catalogue label reports the announced mathematical claim.

\subsection{Sources}
\begin{itemize}
\item OpenAI, \emph{Nonuniqueness for the periodic hard-sphere Boltzmann equation}, 2026-09-23; citation key \texttt{OAI:Nonuniqueness-for-the-periodic-hard-sphere-Boltzmann-equation-September-23-2026}. \href{https://github.com/openai/math/blob/adc7f1241b42e322a6451854ab7e4b4c146bf78a/preprints/Nonuniqueness-for-the-periodic-hard-sphere-Boltzmann-equation-September-23-2026/paper.pdf}{PDF (GitHub)}; \href{https://github.com/openai/math/raw/adc7f1241b42e322a6451854ab7e4b4c146bf78a/preprints/Nonuniqueness-for-the-periodic-hard-sphere-Boltzmann-equation-September-23-2026/paper.pdf}{PDF (direct)}; \href{https://github.com/openai/math/blob/adc7f1241b42e322a6451854ab7e4b4c146bf78a/preprints/Nonuniqueness-for-the-periodic-hard-sphere-Boltzmann-equation-September-23-2026/README.md}{Manuscript page and citation}; \href{https://github.com/openai/math/blob/main/preprints/Nonuniqueness-for-the-periodic-hard-sphere-Boltzmann-equation-September-23-2026/paper.pdf}{Latest version}.
\item OpenAI, \emph{Nonuniqueness with local conservation for the hard-sphere Boltzmann equation}, 2026-10-05; citation key \texttt{OAI:Nonuniqueness-with-local-conservation-for-the-hard-sphere-Boltzmann-equation-October-5-2026}. \href{https://github.com/openai/math/blob/adc7f1241b42e322a6451854ab7e4b4c146bf78a/preprints/Nonuniqueness-with-local-conservation-for-the-hard-sphere-Boltzmann-equation-October-5-2026/paper.pdf}{PDF (GitHub)}; \href{https://github.com/openai/math/raw/adc7f1241b42e322a6451854ab7e4b4c146bf78a/preprints/Nonuniqueness-with-local-conservation-for-the-hard-sphere-Boltzmann-equation-October-5-2026/paper.pdf}{PDF (direct)}; \href{https://github.com/openai/math/blob/adc7f1241b42e322a6451854ab7e4b4c146bf78a/preprints/Nonuniqueness-with-local-conservation-for-the-hard-sphere-Boltzmann-equation-October-5-2026/README.md}{Manuscript page and citation}; \href{https://github.com/openai/math/blob/main/preprints/Nonuniqueness-with-local-conservation-for-the-hard-sphere-Boltzmann-equation-October-5-2026/paper.pdf}{Latest version}.
\item Family 363: \href{https://github.com/openai/math/blob/adc7f1241b42e322a6451854ab7e4b4c146bf78a/lean/docs/363.md}{Lean scope}.
\item \href{https://github.com/openai/math/blob/adc7f1241b42e322a6451854ab7e4b4c146bf78a/lean/ComparatorChallenges/BoltzmannNonuniqueness.json}{Comparator BoltzmannNonuniqueness.json}.
\item \href{https://github.com/openai/math/blob/adc7f1241b42e322a6451854ab7e4b4c146bf78a/lean/ComparatorChallenges/BoltzmannNonuniqueness.lean}{Comparator BoltzmannNonuniqueness.lean}.
\item \href{https://github.com/openai/math/blob/adc7f1241b42e322a6451854ab7e4b4c146bf78a/CONTENTS.md}{OpenAI catalogue entry}.
\end{itemize}
\end{document}
